\section{用随机草图把大矩阵压小}
\label{sec:08-Numerical-Analysis/05-Numerical-Linear-Algebra/08}

设计矩阵有 \(10^7\) 行、\(10^3\) 列：一千万条记录，每条一千个特征，要解的是一个最小二乘。按\hyperref[sec:08-Numerical-Analysis/05-Numerical-Linear-Algebra/02]{``QR 与最小二乘的稳定解法''}那一节的做法走 Householder QR，运算量是 \(O(mn^2)\approx10^{13}\)，把矩阵以双精度整个装进内存要 80 GB，而记录每天还在往里加。

会上有人提议：随机抽一百万行，拿抽出来的子问题当替身。这个提议直觉上很可疑——万一抽掉的正是最关键的那几行呢？

疑虑是对的。但它指向的补救不是``别抽''，而是``别用抽的''。

\subsection{均匀抽行确实可能出事}

构造一个极端情形就够了。设 \(A\) 的第 \(n\) 列只在第一行有一个非零元，其余 \(10^7-1\) 行在这一列上全是零。那么关于 \(x_n\) 的全部信息只写在第一行里。均匀抽一百万行，第一行被抽中的概率是 \(0.1\)；九成的情形下，压缩后的问题在这个坐标上根本无解可定，\(\tilde x_n\) 取什么都行。麻烦的是压缩后的残差看上去毫无问题——它对压缩后的方程组拟合得很好，只是那个方程组已经不再是原来的问题。

这件事有一个标准的度量。记 \(h_i=a_i^{\trans}(A^{\trans}A)^{-1}a_i\) 为第 \(i\) 行的杠杆值，它落在 \([0,1]\) 里，且 \(\sum_i h_i=n\)。杠杆值刻画的是这一行在列空间里有多``独特''：上面那个构造中 \(h_1=1\)，一行独占了总量 \(n\) 里的一整份。均匀抽样只在杠杆值摊得均匀（每行都约为 \(n/m\)）时才安全，而真实数据几乎从不保证这一点。

按杠杆值大小抽样能治好这个病，但绕回了原点：算杠杆值要先有 \((A^{\trans}A)^{-1}\)，也就是先做一次 QR，而 QR 正是我们想避开的那一步。

\subsection{随机投影不做选择，它做混合}

换一个动作。取 \(S\in\R^{s\times m}\)，\(s\ll m\)，用 \(SA\) 与 \(Sb\) 代替 \(A\) 与 \(b\)。这里的关键不在于 \(s\) 小，而在于 \(SA\) 的第 \(j\) 行是 \(\sum_i S_{ji}a_i^{\trans}\)——原矩阵的\emph{每一行}都出现在里面。上一小节那个构造在这里不成立了：第一行没有被丢掉的机会，因为压根没有``挑选''这个动作。

\begin{center}
\begin{tikzpicture}[x=1cm,y=1cm,font=\small,>=stealth]
  \draw[black!50,fill=orange!16] (0.2,3.5) rectangle (3.6,4.4);
  \node[font=\scriptsize,black!80] at (1.9,3.95) {\(S\)};
  \draw[<->,black!45] (0.2,4.7) -- (3.6,4.7);
  \node[above,font=\scriptsize,black!65] at (1.9,4.66) {\(m\)};
  \draw[<->,black!45] (-0.05,3.5) -- (-0.05,4.4);
  \node[left,font=\scriptsize,black!65] at (-0.12,3.95) {\(s\)};
  \node at (3.95,3.95) {\(\times\)};
  \draw[black!50,fill=blue!13] (4.3,0.4) rectangle (5.3,5.4);
  \node[font=\scriptsize,black!80] at (4.8,2.9) {\(A\)};
  \draw[<->,black!45] (4.3,5.65) -- (5.3,5.65);
  \node[above,font=\scriptsize,black!65] at (4.8,5.61) {\(n\)};
  \draw[<->,black!45] (5.6,0.4) -- (5.6,5.4);
  \node[right,font=\scriptsize,black!65] at (5.66,2.9) {\(m\)};
  \node at (6.45,3.95) {\(=\)};
  \draw[black!50,fill=green!18] (6.95,3.5) rectangle (7.95,4.4);
  \node[font=\scriptsize,black!80] at (7.45,3.95) {\(SA\)};
  \node[align=left,anchor=west,font=\footnotesize,black!70] at (8.35,3.95)
    {\(s\) 只由 \(n\) 与 \(1/\epsilon^{2}\) 决定，\\与行数 \(m\) 无关；\\\(m\) 再涨一百倍，\\\(s\) 一个不变};
  \begin{scope}[shift={(0.7,-1.05)}]
    \fill[green!14] (0.7,-0.45) rectangle (9.3,0.45);
    \draw[black!45] (0.7,0.45) -- (9.3,0.45);
    \draw[black!45] (0.7,-0.45) -- (9.3,-0.45);
    \draw[black!45,dashed] (0.7,0) -- (9.3,0);
    \node[left,font=\scriptsize,black!65] at (0.62,0.45) {\(1+\epsilon\)};
    \node[left,font=\scriptsize,black!65] at (0.62,0) {\(1\)};
    \node[left,font=\scriptsize,black!65] at (0.62,-0.45) {\(1-\epsilon\)};
    \draw[thick,blue!62] plot[smooth] coordinates
      {(0.9,0.16) (1.7,-0.21) (2.5,0.28) (3.3,0.04) (4.1,-0.29) (4.9,0.21)
       (5.7,-0.09) (6.5,0.26) (7.3,-0.24) (8.1,0.11) (8.9,-0.17) (9.2,0.14)};
    \node[align=center,font=\footnotesize,black!70] at (5.0,-1.0)
      {比值 \(\norm{SAx}_2/\norm{Ax}_2\)，\(x\) 跑遍整个 \(\R^{n}\)};
  \end{scope}
\end{tikzpicture}

\emph{上排：一个瘦高的矩阵被一个矮胖的随机矩阵从左边压成一个小矩阵，而压完之后的行数只跟列数和精度要求有关。下排：这次压缩要兑现的承诺不是``某几个向量的长度大致不变''，而是列空间里\emph{所有}向量的长度同时被夹在一条窄带内。}
\end{center}

图下排那条窄带就是要证的东西，把它写成定义。

\begin{definition}[子空间嵌入]
称 \(S\in\R^{s\times m}\) 是子空间 \(\mathcal{L}\subseteq\R^m\) 的 \((1\pm\epsilon)\) 嵌入，如果
\[
(1-\epsilon)\norm{y}_2\;\le\;\norm{Sy}_2\;\le\;(1+\epsilon)\norm{y}_2
\qquad\text{对所有 }y\in\mathcal{L}.
\]
\end{definition}

取 \(\mathcal{L}=\rangesp(A)\)，这个条件就是

\[
(1-\epsilon)\norm{Ax}_2\;\le\;\norm{SAx}_2\;\le\;(1+\epsilon)\norm{Ax}_2
\qquad\text{对所有 }x\in\R^n .
\]

\begin{theorem}[Gauss 草图给出子空间嵌入]
设 \(S\) 的元素独立取自 \(N(0,1/s)\)，\(\mathcal{L}\subseteq\R^m\) 是一个事先固定的 \(n\) 维子空间。存在绝对常数 \(C\)，使得只要
\[
s\;\ge\;C\,\frac{n+\log(1/\delta)}{\epsilon^{2}},
\]
则 \(S\) 以概率至少 \(1-\delta\) 是 \(\mathcal{L}\) 的 \((1\pm\epsilon)\) 嵌入。
\end{theorem}

论证骨架分三步。第一步是单个向量：对固定的单位向量 \(y\in\mathcal{L}\)，\(\norm{Sy}_2^2\) 服从 \(\chi^2_s/s\)，标准的次指数集中给出
\[
P\bigl(\,\bigl\lvert\norm{Sy}_2-1\bigr\rvert>\tfrac{\epsilon}{4}\,\bigr)\;\le\;2\exp(-c\,s\,\epsilon^{2}),
\]
这正是 Johnson--Lindenstrauss 引理的核心不等式，\hyperref[sec:13-Machine-Learning/04-Dimensionality-Reduction/05]{``降维与升维的对偶：在合适的维度做合适的事''}把它在点集上的用法讲过一遍。

第二步是从一个向量推到无穷多个。\(\mathcal{L}\) 的单位球面上有连续统那么多向量，逐个套集中界并做并集界是行不通的。补救是取一张 \(\epsilon/8\) 网 \(\mathcal{N}\)：\(n\) 维单位球面上存在这样一张网，元素个数不超过 \((24/\epsilon)^n\)。对网上的每一点用集中界再做并集界，失败概率不超过 \((24/\epsilon)^n\cdot2\exp(-cs\epsilon^2)\)，要它小于 \(\delta\)，取对数就得到定理里 \(s\) 的下界。

第三步把网上的结论抬到整个球面：任取单位向量 \(y\)，找网上最近的 \(y'\)，则 \(\norm{y-y'}_2\le\epsilon/8\)，而 \(\norm{S(y-y')}_2\) 又可以用球面上的上确界自己来控制，移项后上确界被自己的一个小倍数吸收掉，常数因子归并进 \(C\)。

整个论证里 \(m\) 一次都没有出现。第一步的集中界只看 \(s\)，第二步的网只看 \(n\)——这就是行数不进入结论的全部原因，也是本篇最值得记住的一句话。

\subsection{压缩之后解出来的东西差多远}

设 \(\tilde x=\argmin_x\norm{S(Ax-b)}_2\)，\(x_\star\) 是原问题的最优解。把嵌入性质用在 \(A\) 的列与 \(b\) 张成的那个 \(n+1\) 维子空间上（残差向量落在里面），于是
\[
\begin{aligned}
\norm{A\tilde x-b}_2
&\le\frac{1}{1-\epsilon}\,\norm{S(A\tilde x-b)}_2
 \le\frac{1}{1-\epsilon}\,\norm{S(Ax_\star-b)}_2\\
&\le\frac{1+\epsilon}{1-\epsilon}\,\norm{Ax_\star-b}_2 .
\end{aligned}
\]
中间那步用的是 \(\tilde x\) 在压缩问题上最优。\(\epsilon\le1/3\) 时右端不超过 \((1+3\epsilon)\)，于是：在一个只有 \(s\) 行的问题上求解，得到的残差与原问题最优残差只差常数倍。

要留神这句话保的是什么。它保的是残差，不是解。\(A\) 病态时，残差相差 \((1+\epsilon)\) 倍的两个点在参数空间里可以离得很远——\hyperref[sec:08-Numerical-Analysis/01-Floating-Point-and-Error/03]{``条件数衡量问题本身的敏感性''}讲的就是这个落差，它不因为换了随机算法而消失。要的是预测值，草图够用；要的是系数本身并且打算解释它，草图给的保证不够。

\subsection{条件对账}

\textbf{\(s\) 与 \(m\) 无关，这条是核心，也只到此为止。}行数涨十倍，\(s\) 一个不变，这是随机草图值钱的地方。但省下来的不自动是时间：稠密 Gauss 矩阵乘 \(A\) 本身要 \(O(smn)\) 次运算，比直接做 QR 的 \(O(mn^2)\) 还多——\(s\) 通常比 \(n\) 大。真正让方法跑得动的是结构化的 \(S\)。CountSketch 每一列只有一个非零元，位置随机、符号随机，\(SA\) 相当于把每一行随机扔进 \(s\) 个桶之一再带符号累加，一遍数据就够，运算量是 \(O(\operatorname{nnz}(A))\)，与 \(s\) 无关；换来的是 \(s\) 的要求从 \(n/\epsilon^2\) 抬到 \(n^2/\epsilon^2\)。子采样 Hadamard 变换在两者之间：\(s\) 只多一个 \(\log n\) 因子，乘法用快速变换在 \(O(mn\log m)\) 内完成。三种 \(S\) 在``\(s\) 要多大''和``乘一次要多久''这两件事上分别站在不同位置，选哪个取决于 \(A\) 是稠密还是稀疏。

\textbf{对 \(S\) 的分布有实质要求。}集中界依赖元素的次高斯性；Gauss 项、Rademacher 项都满足。若元素取自方差存在但尾巴很重的分布，单个向量的集中界就退化成多项式速率，第二步的并集界再也压不住，定理不成立。CountSketch 不是次高斯的，它的分析走的是另一条路（有界项加上矩方法），结论形式相同而常数不同——这不是同一个定理的推广，是另一个定理。

\textbf{子空间必须在抽 \(S\) 之前就定下来。}定理断言的是``对固定的 \(\mathcal{L}\)，随机的 \(S\) 以高概率好''，不是``对固定的 \(S\)，一切 \(\mathcal{L}\) 都好''。所以同一个 \(S\) 不能反复用在由上一轮结果决定的新子空间上：那时子空间与 \(S\) 相关，概率论证的前提破了。迭代算法里这一条最容易被无声地违反，正确做法是每轮重抽，或者一开始就按最终要覆盖的那个子空间取足够大的 \(s\)。

\textbf{精度以 \(1/\epsilon^2\) 计价，所以草图不是用来求高精度解的。}把 \(\epsilon\) 从 \(10^{-1}\) 压到 \(10^{-3}\)，\(s\) 要涨一万倍，小矩阵早就不小了。成熟的用法是让随机性只负责``快''，让确定性负责``准''：先算 \(SA\) 的 QR 分解 \(SA=QR\)，再把 \(R^{-1}\) 当预条件子——嵌入性质保证 \(AR^{-1}\) 的条件数以高概率贴近 1，于是在\emph{原}问题上跑 LSQR，几十次迭代就到机器精度。此时随机性影响的只是迭代步数，最终精度由确定性的迭代法兜底。这个分工模式值得单独记住：把随机化放在决定速度的位置，把确定性放在决定精度的位置。

\textbf{结论只以概率 \(1-\delta\) 成立。}这与前面七篇给出的全部结论都不是同一种承诺。好在 \(\delta\) 以 \(\log(1/\delta)\) 的形式进入 \(s\) 的下界，把失败概率从 \(10^{-2}\) 压到 \(10^{-6}\) 只让 \(s\) 变成三倍左右。这一条留到本章最后一篇专门讨论。

最后拦一个越界的推论。草图跑通了、下游指标没掉，不能反推``这份数据其实只有 \(s\) 维''。嵌入定理对任何 \(n\) 维子空间都成立，不需要数据本身有低维结构，它保住的只是这一次任务用到的那些长度——\hyperref[sec:14-Deep-Learning/10-Interpretability-Mathematics/03]{``高维数据的低维结构检验''}把这个混淆拆开过，\hyperref[sec:11-Mathematical-Statistics/09-Random-Matrix-Theory/03]{``谱偏差怎样影响 PCA 与随机投影''}则从随机矩阵谱的一侧说明了为什么随机投影的保证不随环境维数恶化。压缩成功只说明当前任务对这次压缩稳健。

本篇的 \(S\) 从左边乘，压的是行数，保的是列空间上的长度。下一篇把随机矩阵挪到右边乘，目标也换一个：不再是把方程变少，而是要探出矩阵最主要的那几个方向。
